\section{First variation of the delay--flux mass estimator}
\label{app:estimator-variation}

The point-mass delay relation is standard; see, for example,
Mu\~noz et al.~\cite{Munoz2016}. This appendix derives its first variation
under a general local lens-potential perturbation, including the displacement
of the images in the magnification calculation. The result is a thin-lens
projection rule; it selects neither a horizon response nor a physical
surface-state amplitude.

Use the reference point-mass Einstein angle and observed time scale
\begin{align}
 \theta_E^2=\frac{4GM}{c^2}\frac{D_{ls}}{D_lD_s},
 \qquad t_M=\frac{4GM}{c^3}(1+z_l),\notag\\
 \qquad \boldsymbol x=\boldsymbol\theta/\theta_E,
 \quad \boldsymbol y=\boldsymbol\beta/\theta_E.
 \label{eq:estimator-scales}
\end{align}
Hold the physical source, body mass, distances, redshift and screen fixed
while varying the dimensionless potential,
\begin{equation}
 \tau_\epsilon(\boldsymbol x;\boldsymbol y)
 =\tfrac12|\boldsymbol x-\boldsymbol y|^2-\ln r
  -\epsilon p(\boldsymbol x),\qquad r=|\boldsymbol x|.
 \label{eq:estimator-fermat}
\end{equation}
For $u=|\boldsymbol y|>0$, put $s=\sqrt{u^2+4}$. On the source axis the
reference images are $x_+=(u+s)/2>1$ and $x_-=-1/x_+<0$. Continue the
positive-parity minimum and negative-parity saddle without crossing a
critical curve or reversing their flux/time ordering. With
$R=|\mu_+|/|\mu_-|>1$ and $\Delta t=t_M(\tau_--\tau_+)>0$,
\begin{align}
 \ln R_0=4\operatorname{asinh}(u/2),\qquad
 h(u)=\tfrac12us+2\operatorname{asinh}(u/2),\notag\\
 M_{\rm app}=\frac{c^3\Delta t}
 {4G(1+z_l)h[2\sinh(\tfrac14\ln R)]}.
 \label{eq:estimator-standard}
\end{align}
Indeed, the lens equation $x-u-1/x=0$ gives
$\mu_+=x_+^4/(x_+^4-1)$, $\mu_-=-1/(x_+^4-1)$ and
$\tau_--\tau_+=\tfrac12(x_+^2-x_+^{-2})+2\ln x_+=h(u)$.
The distance factors in the physical time prefactor cancel against
$\theta_E^2$, so $M_{\rm app}=M$ for every reference alignment. The
aligned Einstein-ring limit is excluded from the two-image estimator.

Let $\delta=\partial_\epsilon|_0$ and define
\begin{equation}
 H_0(\boldsymbol x)=D^2\ln r
 =\frac{I}{r^2}-\frac{2\boldsymbol x\boldsymbol x^T}{r^4},
 \qquad A_i=I-H_0(\boldsymbol x_i).
\end{equation}
Differentiating the stationary lens equation at fixed $\boldsymbol y$ yields
\begin{align}
 \delta\boldsymbol x_i=A_i^{-1}\nabla p(\boldsymbol x_i),\qquad
 \delta A_i=-D^2p(\boldsymbol x_i)
             -DH_0(\boldsymbol x_i)[\delta\boldsymbol x_i],\notag\\
 \delta\ln|\mu_i|=-\operatorname{tr}(A_i^{-1}\delta A_i).
 \label{eq:estimator-total-hessian}
\end{align}
The evaluation-point term is essential. Explicitly,
\begin{align}
 DH_0(\boldsymbol x)[\boldsymbol v]&
 =-\frac{2(\boldsymbol x\cdot\boldsymbol v)I}{r^4}
 -\frac{2(\boldsymbol v\boldsymbol x^T+\boldsymbol x\boldsymbol v^T)}{r^4}\notag\\
 &\quad+\frac{8(\boldsymbol x\cdot\boldsymbol v)
                 \boldsymbol x\boldsymbol x^T}{r^6}.
\end{align}
Stationarity removes image displacements from the first arrival-time
variation, giving $\delta(\tau_--\tau_+)=p(\boldsymbol x_+)-p(\boldsymbol x_-)$.
Since $h'=s$ and $(\ln R_0)'=4/s$, the total estimator variation is
\begin{equation}
 \boxed{\displaystyle
 \delta\ln M_{\rm app}
 =\frac{p(\boldsymbol x_+)-p(\boldsymbol x_-)}{h(u)}
 -\frac{s^2}{4h(u)}
 \bigl[\delta\ln|\mu_+|-\delta\ln|\mu_-|\bigr].}
 \label{eq:estimator-general-variation}
\end{equation}
Thus a pure alignment displacement cancels. A changing physical mass or
redshift prefactor is a separate variation; it must not be hidden in an
unspecified potential normalization.

For the quadrupole $p=q(r)\cos[2(\varphi-\varphi_q)]$, the two reference
image angles differ by $\pi$. Write
$c_2=\cos[2(\varphi_s-\varphi_q)]$, $r_+=(u+s)/2$, $r_-=1/r_+$,
$a_r=1+r^{-2}$ and $a_t=1-r^{-2}$. The radial image displacement is
$\delta x_r=q'c_2/a_r$. Equation~\eqref{eq:estimator-total-hessian} then gives
\begin{align}
 \delta\ln|\mu_i|&=c_2\mathcal F_q(r_i),\\
 \mathcal F_q(r)&=
 \frac{q''}{a_r}+\frac{q'/r-4q/r^2}{a_t}
 +\frac{2q'}{r^3a_r}\left(\frac1{a_r}-\frac1{a_t}\right),
 \label{eq:estimator-Fq}\\
 \delta\ln M_{\rm app}&=c_2\mathcal K_q(u),\\
 \mathcal K_q(u)&=
 \frac{q(r_+)-q(r_-)}{h(u)}
 -\frac{s^2}{4h(u)}\bigl[\mathcal F_q(r_+)-\mathcal F_q(r_-)\bigr].
 \label{eq:estimator-Kq}
\end{align}
The final term of $\mathcal F_q$ is the radial evaluation-point contribution.
Tangential displacements contribute off-diagonal terms whose trace vanishes
in the reference radial/tangential eigenbasis. The physical action or
worldtube source must supply $q,q',q''$ in the same lens units before this
linear projection is applied.

The following closed forms provide nuisance comparisons and independent
checks. Defining $K_*=s^3/[u h(u)]$, direct substitution yields
\begin{equation}
 \mathcal K_{r^2}=K_*,\qquad
 \mathcal K_{r^{-2}}=(u^2+1)K_*,\qquad
 \mathcal K_1=K_*.
 \label{eq:estimator-power-responses}
\end{equation}
Hence a compact projected quadrupole
$B r^{-2}\cos[2(\varphi-\varphi_B)]$ contributes
$B(u^2+1)K_*\cos[2(\varphi_s-\varphi_B)]$, while an external shear in the
convention $(\gamma/2)r^2\cos[2(\varphi-\varphi_\gamma)]$ contributes
$(\gamma/2)K_*\cos[2(\varphi_s-\varphi_\gamma)]$. At one alignment the
amplitudes are degenerate; several alignments distinguish these two
specified shapes only to the extent that the nuisance assumptions hold.

Linearity also gives a nontrivial blind direction,
\begin{equation}
 \boxed{\mathcal K_{r^2-1}(u)=0\qquad\text{for every }u>0.}
 \label{eq:estimator-blind-direction}
\end{equation}
This is an exact \emph{first-order} null of the statistic, not an exact
finite-perturbation equivalence or a vanishing lens deformation. The potential
$(r^2-1)\cos2\varphi$ generally changes images, delay and flux ratio.
It is a local smooth-potential counterexample on image annuli; whether a
specified physical source can generate it requires its own source/boundary
analysis. Different radial potentials therefore do not automatically give
identifiable tracker signals. An inference must compare the evaluated
$\mathcal K_q$ with the justified nuisance response space, including its
conditioning, rather than compare potential labels.

These formulas concern quasistatic geometric optics and two noncritical
continued primary images. They do not describe near-horizon strong-deflection
images, new-image total flux, differential extinction or unmatched variable
source emission. Exact symbolic identities, Cartesian total-Hessian checks
and finite on-axis polynomial branch checks are supplied with the
reproduction sources; they are internal consistency tests, not observations.
